% SOURCE_AUTHORITY: SGA5_LNM589_SCAN, PDF pages 5--7 = printed Roman pages V--VII (scan-only authority).
% SOURCE_DOCUMENT_SHA256: B256EBD072A8C68209518412A263C9289C6F1854A346733D86F885930D5FE6CA
Após as primeiras análises técnicas da exposição V sobre sistemas projetivos ádicos, a exposição VI define a noção de $\mathbb Z_\ell$-feixe construtível (resp., de feixe localmente constante construtível, isto é, liso na terminologia de Deligne ([1], (SGA~$4^{1/2}$ Arcata)) e estende aos $\mathbb Z_\ell$-feixes construtíveis, por passagem ao limite a partir dos $\mathbb Z/\ell^n$-feixes, o formalismo das imagens diretas superiores com suportes próprios. A exposição VI, contudo, não define uma categoria derivada dos $\mathbb Z_\ell$-feixes, munida das operações $\lten$ e $\uRHom$ que prolongariam as disponíveis na categoria derivada dos $\mathbb Z/\ell^n$-feixes. \emph{A fortiori}, tampouco se considera a extensão aos $\mathbb Z_\ell$-feixes do formalismo de dualidade de SGA~4 XVIII. A questão foi objeto da tese (não publicada) de Jouanolou. O leitor poderá consultar, entretanto, o início de [1], onde se definem os $\uExt^i$ dos $\mathbb Z_\ell$-feixes.

A exposição VII, independente do restante do seminário, é sem dúvida uma das mais interessantes. Contém: a) o cálculo da cohomologia de algumas variedades clássicas (espaços afins, espaços projetivos e variedades de bandeiras); b) uma exposição da teoria das classes de Chern em cohomologia étale; c) uma demonstração da fórmula de autointerseção $i^*i_*x=x\,c_d(N)$ em cohomologia étale (respectivamente, no anel de Chow, demonstração devida nesse caso a Mumford), bem como várias aplicações dessa fórmula (cálculo da cohomologia das variedades obtidas por explosão e fórmula de Gauss--Bonnet).

% SOURCE_EMENDATION: printed Roman p. V repeats «L'exposé VII» here. The subject (finiteness conditions and Grothendieck groups), the printed contents, and the dependency diagram identify this paragraph as Exposé VIII; Spanish corrects the numeral accordingly.
A exposição VIII contém sorites sobre as condições de finitude nas categorias derivadas e nos grupos de Grothendieck. Essas questões são retomadas em um quadro mais geral em (SGA~6 I, IV).

A exposição IX, de J.-P. Serre, intitulada Introdução à teoria de Brauer, foi publicada em [4] e não é reproduzida neste volume. Um de seus principais resultados, devido a Swan e relativo à existência de um módulo projetivo cujo caráter é o caráter de Swan, é usado na exposição X, onde se demonstra uma fórmula de Euler--Poincaré para um feixe sobre uma curva, ligeiramente mais geral do que a que figura na exposição de Raynaud [3].

% SOURCE_PAGE_BOUNDARY: end of printed Roman p. V.
\clearpage

A exposição XI, redigida por I. Bucur, não existe, pois se perdeu durante uma mudança e o autor não havia guardado nenhuma cópia. Era dedicada à teoria de Grothendieck dos traços comutativos, que generaliza a de Stallings [5]. Uma versão mais elaborada dessa teoria, em um quadro de teoria de feixes, é exposta em (III~B 5, 6). As exposições XII e III~B apresentam aplicações da fórmula de Lefschetz. A fórmula de traço da exposição XII é demonstrada independentemente da fórmula geral da exposição III, mas em (III~B 6) mostra-se que os termos locais que nela figuram são precisamente os da fórmula geral e que esta última a implica.

A ausência da exposição XIII é apenas por uma razão de numeração: inicialmente, o conteúdo das exposições X a XII foi previsto para ser expandido para quatro exposições.

A exposição XIV (às vezes também numerada como XV) contém generalidades sobre a correspondência de Frobenius (entre elas, a comparação entre o Frobenius geométrico e o Frobenius aritmético) e, a partir da fórmula de traço da exposição XII (ou III~B), apresenta a demonstração de Grothendieck da racionalidade das funções $L$: os dois pontos importantes são a redução ao caso das curvas e a passagem ao limite, que permite reduzir o problema a uma fórmula para $\mathbb Z/\ell^n$-feixes.

Durante o seminário oral, Grothendieck proferiu uma exposição sobre problemas em aberto e enunciou algumas conjecturas. Essa exposição, que deveria encerrar o seminário, infelizmente não foi redigida, tampouco o foi sua bela exposição introdutória, na qual passou em revista as fórmulas de Euler--Poincaré e de Lefschetz em vários contextos (topológico, analítico complexo e algébrico).

Agradeço a P. Deligne por ter me convencido a escrever, em uma nova versão da exposição III, uma demonstração da fórmula de Lefschetz--Verdier, eliminando assim um dos obstáculos à publicação deste seminário. Sou-lhe muito grato pelas melhorias de redação que me sugeriu,

% SOURCE_PAGE_BOUNDARY: the printed sentence continues from Roman p. VI to p. VII here.
\clearpage
\noindent e pela ajuda que me prestou na preparação deste volume para a editora. Agradeço também à Sra.~Lièvremont, que datilografou com cuidado, e em pouco tempo, as exposições III e III~B.

Não posso terminar esta introdução sem homenagear a memória de I. Bucur, que faleceu de câncer em setembro de 1976. Quem, como eu, teve o privilégio de tê-lo como amigo não esquece sua delicadeza requintada, seu humor sorridente e o encanto de sua conversa.

\begin{flushright}
Paris, 19 de fevereiro de 1977\\[1.2em]
Luc Illusie
\end{flushright}
